\makeabstract
    {
    The Effect of Outer Giant Planets on Inner System Architectures
    }
    {
    Simone Lilavois\textsuperscript{1}, Tiger Lu\textsuperscript{2}, Phil Armitage\textsuperscript{2}, Daniella Bardalez Gagliuffi\textsuperscript{1}
    }
    {
    \textsuperscript{1} Amherst College, Amherst, MA\\
\textsuperscript{2} Center for Computational Astrophysics, Flatiron Institute, NY
    }
    {
    Understanding the processes that dictate planetary formation and evolution is crucial to constraining how the architectures of exoplanetary systems emerge, as well as for determining regimes of stability and chaos. Observations from NASA{\textquoteright}s Kepler mission revealed a remarkable trend of intra-system uniformity in inner terrestrial planet masses, radii, and orbital spacings, a phenomenon referred to as {\textquotedblleft}peas in a pod{\textquotedblright} architectures. Evidence suggests that these system architectures may result from later dynamical sculpting, rather than being a natural outcome of early planet formation. Furthermore, systems featuring an orbital gap between two of the inner compact planets preferentially harbor an outer gas giant, suggesting a possible dynamical process whereby the outer giant induces instabilities, collisions, and ejections that reshape inner system architectures. 
    }
    {
    Here, we investigate whether the inward migration of an outer giant planet can sweep strong enough gravitational resonances to adequately pump the eccentricities and inclinations of inner planets to generate the observed gaps. Using the N-body integrator REBOUND, we simulate randomized inner compact terrestrial systems and track their evolution as a Jupiter-mass planet migrates inward. We complement these numerical simulations with analytical Laplace-Lagrange secular theory calculations using celmech to evaluate mean motion resonances and secular resonance crossings over time.
    }
    {
    Preliminary results demonstrate that a single outer giant is too weakly coupled to the terrestrial inner planets to sweep strong secular resonances. Moreover, significant differences in semi-major axes between the outermost terrestrial planet and the giant planet ensure that only high-order mean motion resonances are encountered during migration, which are too weak to trigger dynamical sculpting. 
    }
    {
    Thus, our early results suggest that a second giant planet may be required to sweep resonances more strongly, or that another unaccounted-for dynamical process is shaping the system architectures we observe. 
    }

    \newpage
    